\section{Security and Recovery Scope}
\label{sec:security}

\subsection{Threat model}
\label{sec:case_study}

AEGIS-Q considers an attacker who can read, copy, modify, and replay the backing directory but cannot obtain the mount credential or compromise the running kernel, FUSE daemon, CUDA driver, or cryptographic libraries.  Under that scope, AES-GCM protects the confidentiality and authenticity of published block records, while the HMAC-protected key envelope and checkpoint detect unauthenticated metadata changes.  The adversary can still deny service, delete records, replay a complete file-backed snapshot, or exploit vulnerabilities outside the FUSE process.  We make no claim against a privileged local attacker that can read process memory, substitute the CUDA driver, alter the mount executable, or use GPU microarchitectural side channels.

The accepted encrypted state consists of a validated DEK envelope, a complete journal mapping, and ciphertext records referenced by that mapping.  A tag failure causes a read to fail.  A malformed or unauthenticated envelope/checkpoint causes open to fail.  These behaviors are necessary for a safe FUSE format, but they are not a formal proof of crash consistency.  In particular, xattr atomicity and directory durability are delegated to the underlying filesystem.

\subsection{Freshness is conditional on an external anchor}
\label{sec:case_problem}

The source tree contains two anchor modes.  The file mode is only an integrity witness; copying the directory copies the witness.  The corrected crash/replay run therefore reports four of four file-backed snapshot restores as \emph{rollback visible}, not as successful protection.  This negative result is important: an HMAC alone does not create freshness when its entire state is replayable.

The hardware mode stores an authenticated prefix record in an administrator-provisioned TPM NV index.  It is deliberately outside the evaluated configuration.  The paper does not assert that an NV index is a signer, does not report an invented TPM latency, and does not claim PCR binding.  To make this mode publishable, future work must document the TPM model/bus, authorization and index attributes, monotonic update protocol, software-update recovery, and a power-fail matrix that distinguishes acknowledged, anchored, and replayed writes.

\subsection{Side channels and application semantics}
\label{sec:case_impl}

GPU SIMT execution, cache behavior, memory coalescing, occupancy, and contention can all leak information.  The repository contained generators that produced simulated timing and TVLA traces.  Those files are excluded from this paper.  The retained GPU integrity test checks output parity with OpenSSL; it says nothing about constant-time execution, physical power/EM leakage, or cross-tenant GPU channels.

Likewise, the clean-remount test validates one encrypted 128~KiB round trip through the final binary.  It does not certify SQLite WAL, \texttt{mmap}, locking, crash recovery across every cut point, or the complete POSIX contract.  We retain the following labels to make the scope explicit: the prior SQLite case study (Section~\ref{sec:case_results}) and edge/cloud comparison (Section~\ref{sec:case_edge_cloud}) are not claimed by this revision.

\subsection{What the current tests establish}
\label{sec:case_results}

The final validation runs a FUSE self-test, clean encrypted remount, a managed-memory smoke test, and GPU integrity parity checks.  The round-trip test confirms that after a clean unmount/remount, the original 128~KiB plaintext is recovered and the corresponding \texttt{.pqcdata} sidecar does not contain the whole plaintext byte sequence.  This is evidence of the repaired key-envelope path, not a statistical security proof.

\subsection{Deployment boundary}
\label{sec:case_edge_cloud}

AEGIS-Q is a prototype for local edge storage.  It makes no edge-versus-cloud cost, latency, or privacy quantification.  The appropriate comparison is deployment-specific and requires a separately measured network/service configuration.  This restriction avoids turning a local microbenchmark into a claim about cloud databases or remote key-management systems.
